
%%%%%%%%%%%%%%%%%%%%%%%%%%%%%%%%%%%%%%%%%%%%%%%%%%%%%%%%%%%%%
\subsection{问题四的模型的建立和求解}
\subsubsection{具体分析}

问题四在问题三的基础上引入定向干扰源，其余条件完全相同，属于不完全观测条件下的序贯决策问题，且其不完全程度较问题三显著加深。问题的质变源于定向源的空间辐射特性：定向源仅在以其定向方向为中心、两侧各 $90^\circ$ 的范围内辐射信号，一旦检测点落入其背向半平面，即便距离远小于有效接收半径也收不到任何信号。这使得"无信号"这一反馈从问题三中的弱信息退化为完全的三义事件，其可能原因包括该频道空闲、检测点超出有效接收半径以及检测点位于背向半平面。若仍沿用问题三的判据，机器狗会把大量实际存在的定向源误判为空闲频道而漏清除，直接违反"确保所有干扰源被清除"的硬约束。因此本问的核心任务是在定向方向未知且接口不返回任何朝向信息的前提下，设计一种使漏测概率为零的观测机制，并尽可能压低为此付出的额外时间代价。

\subsubsection{模型建立}

\paragraph{可检测域的半圆盘刻画} 设定向源位于 $G_c$、有效接收半径为 $r_c$、定向方向为 $\psi_c$，则检测点 $P$ 可被检测的充要条件为
\begin{equation}
\lVert P-G_c\rVert\le r_c
\quad\text{且}\quad
\bigl\lvert\operatorname{wrap}\bigl(\arg(P-G_c)-\psi_c\bigr)\bigr\rvert\le90^\circ ,
\label{eq:directional}
\end{equation}
其中 $\operatorname{wrap}(\cdot)$ 把角度归化到 $[-180^\circ,180^\circ)$。式 \eqref{eq:directional} 表明定向源的可检测域是圆盘与以 $\psi_c$ 为外法向的闭半平面的交，即一个半圆盘；而全向源的可检测域为整个圆盘。这一几何差异是问题四全部算法的出发点。

\paragraph{背向不可测定理} 若两个探测点 $P_1,P_2$ 对 $G_c$ 的视角差不少于 $180^\circ$ 且二者到 $G_c$ 的距离都不超过 $r_c$，则二者不可能同时返回无信号。

该定理的证明是直接的：由式 \eqref{eq:directional}，若 $P_i$ 返回无信号且 $\lVert P_i-G_c\rVert\le r_c$，则 $P_i$ 必落在背向半平面 $\{X:(X-G_c)\cdot u(\psi_c)<0\}$ 内；半平面的张角为 $180^\circ$，从 $G_c$ 观察落入同一开半平面内的两点的方位角之差必然严格小于 $180^\circ$，与视角差不小于 $180^\circ$ 矛盾。该定理给出一条可操作的判别准则：若某一频道在两个视角差接近或超过 $180^\circ$ 的观测点处均无信号，且两点到候选位置的距离都在有效接收半径内，则该频道必然空闲，可以从待办列表中安全移除；反之若两个无信号点的视角差很小，则不能排除定向遮蔽，必须继续补测。

\paragraph{方位包围判据} 定理的逆否形式给出了更一般的要求。设机器狗在位置集合 $\mathcal Q$ 上观测，若频道 $c$ 全程返回无信号，则真实源 $G_c$ 必满足所有落于 $\mathcal B(G_c,r_c)$ 内的观测点都位于背向半平面。由于定向方向 $\psi_c$ 未知，要保证不漏测，必须对所有可能的 $\psi$ 都不存在这样的情形，即
\begin{equation}
\forall\psi,\ \exists P\in\mathcal Q:\ \lVert P-G_c\rVert\le r_c\ \text{且}\ (P-G_c)\cdot u(\psi)\ge0 .
\label{eq:zero-miss}
\end{equation}
式 \eqref{eq:zero-miss} 等价于一个纯几何条件
\begin{equation}
\boxed{\ \text{零漏测}\iff
\mathcal Q\cap\mathcal B(G_c,r_c)\ \text{的方位角最大间隔}<180^\circ
\iff G_c\in\operatorname{int}\operatorname{conv}\bigl(\mathcal Q\cap\mathcal B(G_c,r_c)\bigr).\ }
\label{eq:criterion}
\end{equation}
即机器狗在干扰源有效接收范围内的观测点必须在方位上"包围"该干扰源。作为特例，若这些点的方位角互差 $120^\circ$，则三点的最小包围弧长为 $240^\circ>180^\circ$，故至少一点必落在覆盖半平面内。

式 \eqref{eq:criterion} 的工程含义十分明确：零漏测要求"包围"而非仅"覆盖"，这解释了为什么问题三的保证覆盖网在定向场景下必然失效——覆盖只保证距离，不保证方位上的包围。

\paragraph{分级冗余网模型} 由式 \eqref{eq:criterion}，冗余网的设计问题可表述为
\begin{equation}
\min\ \lvert\mathcal Q\rvert
\quad\text{s.t.}\quad
\max_{\text{方位间隔}}\bigl(\mathcal Q\cap\mathcal B(G,r_c)\bigr)<180^\circ,
\quad\forall G\in D .
\label{eq:net-design}
\end{equation}
本文采用"中心加分级同心环"的构型族求解式 \eqref{eq:net-design}。采用分级而非等半径的理由是：靠近圆心的干扰源只需少量近点即可被包围，而靠近边界的干扰源其可观测点全部集中在朝向圆心的一侧，必须在圆域之外追加观测点才能形成包围。对外环半径 $\rho_o$ 与边界源可见窗口的关系可由余弦定理解析给出
\begin{equation}
\lVert P-G\rVert\le r_c
\iff
\cos\theta\ \ge\ \frac{\rho_o^2+R^2-r_c^2}{2\rho_o R},
\label{eq:window}
\end{equation}
其中 $\theta$ 为观测点在极坐标下的方位角与源方位的夹角。式 \eqref{eq:window} 说明可见窗口的半角随外环半径增大而单调收窄，因此外环不能放得太远。

\paragraph{自适应补测策略} 静态冗余网虽然能给出确定性的漏测保证，但其代价是观测点数成倍增长。考虑到实际场景中定向源只占一部分、且大部分定向源会被常规覆盖网直接检出，本文给出自适应补测策略：先用问题三的七点覆盖网完成一轮全频道粗探；对返回过示向度的频道按问题三流程正常清除；对始终无信号的频道，仅在存在未解频道时才沿加密环（中环与外环）对这些频道做补充观测。由于加密环只需对少数未解频道复测，其检测与切换开销远低于静态方案的全频道扫描。

对于需要定向补测的频道，还可利用历史观测点的方位分布对定向方向作极大似然估计
\begin{equation}
\hat\psi_c=\arg\max_{\psi}\prod_{P\in\mathcal P_c}
\mathbb I\bigl[\lvert\operatorname{wrap}(\arg(P-\hat G_c)-\psi)\rvert>90^\circ\bigr],
\label{eq:psi-mle}
\end{equation}
即在使全部历史观测点都落入背向半平面的方向集合中取中点。该估计在观测点数不少于 $2$ 时可由两两视角差的补集直接给出闭式解，据此机器狗可沿 $\hat\psi_c$ 方向绕行至候选位置另一侧补测一次，若该点与既有观测点的视角差接近 $180^\circ$，则无论结果是信号还是无信号都能给出确定判据。

\subsubsection{模型求解}

对背向不可测定理做数值验证：随机生成干扰源位置与定向方向，并在其 $1000$ 米邻域内构造视角差不小于 $180^\circ$ 的探测点对，统计两点同时返回无信号的次数，结果为 $0$ 组，与定理断言完全一致。作为对照，单点探测的无信号频率为 $0.4960$，与理论值 $0.5$（背向半平面占总方位的 $180/360$）吻合到千分位，说明遮蔽机理的建模正确。

\begin{figure}[htbp]
	\centering
	\includegraphics[width=0.8\textwidth]{../求解/问题四/figures/定向源可检测域与判别定理图.png}
	\caption{定向源的可检测域与两个判别定理的几何示意}
	\label{fig:q4-theory}
\end{figure}

如图 \ref{fig:q4-theory}(a) 所示，全向源与定向源的可检测域分别为完整圆盘与半圆盘，背向半平面内完全没有辐射。图 \ref{fig:q4-theory}(b) 展示背向不可测定理：位于定向方向一侧的探测点返回信号，而对径方向的探测点返回无信号，二者视角差 $180^\circ$。图 \ref{fig:q4-theory}(c) 展示三方位定理：三个互差 $120^\circ$ 的观测点中，位于覆盖半平面内的点返回信号，落入背向半平面的点返回无信号，但不会全部无信号。

对问题三的七点覆盖网按式 \eqref{eq:criterion} 评估方位间隔，结果表明其全域最大方位间隔达到 $360^\circ$，即存在大量位置仅有一个观测点落在有效接收半径内。若定向方向随机均匀分布，则某位置的期望漏测概率等于该处最大方位间隔与 $360^\circ$ 之比，全域平均漏测概率为 $65.24\%$。按距原点半径分层统计的结果如表 \ref{tab:q4-layer} 所示。

\begin{table}[H]
	\centering
	\caption{七点覆盖网对定向源的漏测风险分层统计}
	\label{tab:q4-layer}
	\begin{tabular}{ccccc}
		\toprule
		半径区间 (m) & 抽样点数 & 最大方位间隔 & 平均漏测概率 & 必然漏测点占比 \\
		\midrule
		$0\sim400$     & $10368$ & $360.00^\circ$ & $63.84\%$ & $50.52\%$ \\
		$400\sim800$   & $10080$ & $176.48^\circ$ & $40.19\%$ & $0.00\%$ \\
		$800\sim1200$  & $10089$ & $360.00^\circ$ & $55.03\%$ & $54.62\%$ \\
		$1200\sim1600$ & $10359$ & $360.00^\circ$ & $85.10\%$ & $100.00\%$ \\
		$1600\sim1800$ & $5184$  & $360.00^\circ$ & $96.97\%$ & $100.00\%$ \\
		\bottomrule
	\end{tabular}
\end{table}

表 \ref{tab:q4-layer} 给出了一个结构性结论：只有在半径 $400\sim800$ 米的环带内，七点网的最大方位间隔小于 $180^\circ$、可以保证不漏测；而在半径超过 $1200$ 米的区域，该网络对定向源必然漏测。其原因是外圈相邻观测点的间距达 $1123$ 米，边界处的干扰源往往只落在一个观测点的接收范围内，单个观测点无法在方位上构成包围。这意味着问题三的策略在问题四中不能直接沿用，必须引入方位冗余。

\begin{figure}[htbp]
	\centering
	\includegraphics[width=0.8\textwidth]{../求解/问题四/figures/七点网可隐藏区与冗余网对比图.png}
	\caption{七点网下的可漏测区与分级零漏测冗余网}
	\label{fig:q4-hidden}
\end{figure}

如图 \ref{fig:q4-hidden}(a) 所示，以极坐标热力图呈现七点网的漏测风险分布，颜色以 $180^\circ$ 为中性点：蓝色表示该位置被观测点包围、对任意定向方向都不会漏测，红色表示存在漏测方向。图中仅有一条狭窄的蓝色环带（约在半径 $600$ 米附近），其余区域均为红色，与表 \ref{tab:q4-layer} 的分层结论一致。

在"中心加分级同心环"构型族内求解式 \eqref{eq:net-design}，得到满足零漏测的最小构型为 $37$ 个观测点，其结构为中心点 $1$ 个、半径 $1000$ 米的 $6$ 边形环、半径 $1300$ 米的 $12$ 边形环与半径 $1900$ 米的 $18$ 边形环，全域最大方位间隔为 $155.67^\circ$，严格小于 $180^\circ$ 门槛。作为对照，$30$ 点以内的构型最大方位间隔均在 $200^\circ$ 以上，$25$ 点构型在半径 $1136$ 米处仍有 $195.81^\circ$ 的最大间隔，说明该规模下界不是构造的偶然结果，而是由"边界处必须形成包围"这一几何要求决定的。

\begin{table}[H]
	\centering
	\caption{问题三与问题四观测方案对比}
	\label{tab:q4-cmp}
	\begin{tabularx}{\textwidth}{LCCCCCC}
		\toprule
		方案 & 观测点数 & 最大方位间隔 & 零漏测 & 巡回路径 (m) & 检测次数 & 粗探时间 (s) \\
		\midrule
		问题三·七点保证覆盖网 & $7$ & $360.00^\circ$ & 否 & $6738$ & $140$ & $2180.6$ \\
		问题四·分级零漏测冗余网 & $37$ & $155.67^\circ$ & 是 & $21870$ & $740$ & $8776.9$ \\
		\bottomrule
	\end{tabularx}
\end{table}

由表 \ref{tab:q4-cmp}，静态零漏测方案的粗探时间为 $8776.9$ 秒，是问题三的 $3.025$ 倍，即定向源引入的时间惩罚系数为 $\kappa_{\mathrm{static}}=(T^{(4)}-T^{(3)})/T^{(3)}=3.025$。这一代价主要来自两方面：观测点数由 $7$ 增至 $37$，使检测次数、切换次数与巡回路径同步增长。

\begin{figure}[htbp]
	\centering
	\includegraphics[width=0.8\textwidth]{../求解/问题四/figures/定向源时间惩罚与网规模权衡图.png}
	\caption{定向源引入的时间代价与观测网规模的权衡}
	\label{fig:q4-tradeoff}
\end{figure}

如图 \ref{fig:q4-tradeoff}(a) 所示，三种方案的粗探时间分别为七点网 $2181$ 秒、自适应补测 $6562$ 秒与静态冗余网 $8777$ 秒；图 \ref{fig:q4-tradeoff}(b) 给出观测点数与全域最大方位间隔的关系曲线，曲线在 $37$ 点处穿过 $180^\circ$ 门槛，此后继续加密的边际收益迅速衰减。

自适应策略的推理是：实际场景中定向方向随机，七点网的平均漏测概率为 $65.24\%$，但被漏测的频道只占少数；对这批少数频道沿加密环（中环 $12$ 点、外环 $18$ 点，巡回路径 $20106$ 米）复测，即可将其分辨为空闲或定向。表 \ref{tab:q4-adapt} 给出不同未解频道数下的时间代价。

\begin{table}[H]
	\centering
	\caption{自适应补测策略的时间代价}
	\label{tab:q4-adapt}
	\begin{tabular}{cccccc}
		\toprule
		未解频道数 & $0$ & $1$ & $2$ & $3$ & $6$ \\
		\midrule
		追加检测次数 & $0$ & $30$ & $60$ & $90$ & $180$ \\
		补测时间 (s) & $0$ & $4201.2$ & $4381.2$ & $4561.2$ & $5101.2$ \\
		自适应总时间 (s) & $2180.6$ & $6381.8$ & $6561.8$ & $6741.8$ & $7281.8$ \\
		相对静态方案节省 & $75.2\%$ & $27.3\%$ & $25.2\%$ & $23.2\%$ & $17.0\%$ \\
		\bottomrule
	\end{tabular}
\end{table}

由表 \ref{tab:q4-adapt} 可见，若七点网未能检出任何频道异常，自适应策略的代价与问题三完全相同；即使存在 $2$ 个未解频道，其总时间 $6561.8$ 秒仍比静态方案节省 $25.2\%$；只有在未解频道数达到 $5$ 个以上时，自适应方案才不再占优。因此本文推荐以自适应补测作为实际执行策略，而以静态冗余网作为保证零漏测的代价上界与兜底方案。

按题目要求，问题四三次正式测试的结果如表 \ref{tab:q4-test} 所示。

\begin{table}[H]
	\centering
	\caption{问题 4 正式测试结果}
	\label{tab:q4-test}
	\begin{tabular}{cccc}
		\toprule
		测试案例编码 & 清除干扰源个数 & 平均定位清除时间 (s) & 程序运行时间 (s) \\
		\midrule
		测试 1 案例编码 & & & \\
		测试 2 案例编码 & & & \\
		测试 3 案例编码 & & & \\
		\bottomrule
	\end{tabular}
\end{table}
